
\documentclass{article}
\usepackage{colortbl}
\usepackage{makecell}
\usepackage{multirow}
\usepackage{supertabular}

\begin{document}

\newcounter{utterance}

\centering \large Interaction Transcript for game `cladder', experiment `full\_v1.5\_default', episode 5573 with Qwen3.5{-}9B.
\vspace{24pt}

{ \footnotesize  \setcounter{utterance}{1}
\setlength{\tabcolsep}{0pt}
\begin{supertabular}{c@{$\;$}|p{.15\linewidth}@{}p{.15\linewidth}p{.15\linewidth}p{.15\linewidth}p{.15\linewidth}p{.15\linewidth}}
    \# & \multicolumn{2}{c}{Player} && \multicolumn{2}{c}{Game Master} \\
    \hline

    \theutterance \stepcounter{utterance}  
    & & & \multicolumn{4}{p{0.6\linewidth}}{
        \cellcolor[rgb]{0.9,0.9,0.9}{
            \makecell[{{p{\linewidth}}}]{
                \texttt{\tiny{[P1$\langle$GM]}}
                \texttt{You are an expert in causal inference. The following question is not a typical commonsense query, but rather a meticulously designed question created by a professor specializing in causal inference, intended to assess the students' mastery of the course content.} \\
\texttt{Imagine a self{-}contained, hypothetical world with only the following conditions, and without any unmentioned factors or causal relationships: Gender has a direct effect on treatment and peanut allergy. Treatment has a direct effect on peanut allergy. For patients who are not male and not receiving treatment, the probability of being allergic to peanuts is 98\%. For patients who are not male and receiving treatment, the probability of being allergic to peanuts is 87\%. For patients who are male and not receiving treatment, the probability of being allergic to peanuts is 34\%. For patients who are male and receiving treatment, the probability of being allergic to peanuts is 14\%. The overall probability of male gender is 43\%. For patients receiving treatment, would it be less likely to see being allergic to peanuts if the patient had received no treatment?} \\
\texttt{Start your answer with 'Yes' or 'No', followed by additional reasoning or evidence to support your explanation.} \\
            }
        }
    }
     \\ \\

    \theutterance \stepcounter{utterance}  
    & \multicolumn{4}{p{0.6\linewidth}}{
        \cellcolor[rgb]{0.561,0.737,0.561}{
            \makecell[{{p{\linewidth}}}]{
                \texttt{\tiny{[P1$\rangle$GM]}}
                \texttt{No} \\
\\ 
\texttt{**Reasoning:**} \\
\\ 
\texttt{To answer this question, we must determine the **Conditional Average Treatment Effect (CATE)** for the specific subgroup of patients who are receiving treatment. The question asks: "For patients receiving treatment, would it be less likely to see being allergic to peanuts if the patient had received no treatment?"} \\
\\ 
\texttt{Mathematically, this is asking if \$P(Y=1 | T=1) < P(Y=1 | T=0)\$ within the context of the treated population. In causal inference terms, we are comparing the observed outcome of the treated group against the **counterfactual** outcome of what would have happened to that same group if they had not received treatment (\$Y\_\{0|1\}\$).} \\
\\ 
\texttt{1.  **Identify the Observed Outcome for the Treated Group:**} \\
\texttt{    From the prompt: "For patients who are not male and receiving treatment, the probability... is 87\%. For patients who are male and receiving treatment, the probability... is 14\%."} \\
\texttt{    We need the overall probability of allergy for the treated group (\$P(Y=1 | T=1)\$).} \\
\texttt{    Let \$M\$ be Male and \$T\$ be Treatment.} \\
\texttt{    \$P(M) = 0.43\$, so \$P(\textbackslash{}neg M) = 0.57\$.} \\
\texttt{    \$P(Y=1 | \textbackslash{}neg M, T=1) = 0.87\$} \\
\texttt{    \$P(Y=1 | M, T=1) = 0.14\$} \\
\\ 
\texttt{    \$P(Y=1 | T=1) = P(Y=1 | \textbackslash{}neg M, T=1)P(\textbackslash{}neg M) + P(Y=1 | M, T=1)P(M)\$} \\
\texttt{    \$P(Y=1 | T=1) = (0.87 \textbackslash{}times 0.57) + (0.14 \textbackslash{}times 0.43)\$} \\
\texttt{    \$P(Y=1 | T=1) = 0.4959 + 0.0602 = 0.5561\$ (55.61\%)} \\
\\ 
\texttt{2.  **Identify the Counterfactual Outcome (\$Y\_\{0|1\}\$):**} \\
\texttt{    We need to estimate \$P(Y=1 | T=0, M)\$ for the treated individuals. Since we do not have direct data on the counterfactuals for the treated group, we must rely on the **ignorability assumption** (or conditional exchangeability). The prompt states: "Imagine a self{-}contained, hypothetical world with only the following conditions... without any unmentioned factors." This implies that within the strata defined by Gender (\$M\$), Treatment assignment is independent of potential outcomes.} \\
\\ 
\texttt{    Therefore, for a patient with a specific gender who is currently treated, their counterfactual outcome (what would happen if they} \\
            }
        }
    }
    & & \\ \\

    \theutterance \stepcounter{utterance}  
    & & & \multicolumn{2}{p{0.3\linewidth}}{
        \cellcolor[rgb]{0.9,0.9,0.9}{
            \makecell[{{p{\linewidth}}}]{
                \texttt{\tiny{[GM$|$GM]}}
                \texttt{No} \\
            }
        }
    }
    & & \\ \\

    \theutterance \stepcounter{utterance}  
    & & & \multicolumn{2}{p{0.3\linewidth}}{
        \cellcolor[rgb]{0.9,0.9,0.9}{
            \makecell[{{p{\linewidth}}}]{
                \texttt{\tiny{[GM$|$GM]}}
                \texttt{no} \\
            }
        }
    }
    & & \\ \\

    \theutterance \stepcounter{utterance}  
    & & & \multicolumn{2}{p{0.3\linewidth}}{
        \cellcolor[rgb]{0.9,0.9,0.9}{
            \makecell[{{p{\linewidth}}}]{
                \texttt{\tiny{[GM$|$GM]}}
                \texttt{game\_result = WIN} \\
            }
        }
    }
    & & \\ \\

\end{supertabular}
}

\end{document}
